% answers_short.tex —— 1983 年·全国卷（文科） · 简短答案速查
% 由 tools/make_short_answers.py 自 answers.tex 的 \gksolhead/\gkanswer 提取，
% 并与 1983/1983_数学真题_文_答案要点.md 逐题比对后生成，勿手改（改源 MD 后重跑）。

\begin{gknotice}
\textbf{使用说明}：本节只给最终答案便于快速核对，完整推导见解析版。答案为编写者独立推导，\textbf{非官方标准答案}。
\end{gknotice}

\gkpart{题目}
\begin{itemize}[leftmargin=2.6em]
  \item[\textbf{1.}] D
  \item[\textbf{2.}] A
  \item[\textbf{3.}] D
  \item[\textbf{4.}] C
  \item[\textbf{5.}] C
  \item[\textbf{6.}] $x^2+y^2=2x$ 表示以 $(1,0)$ 为圆心、$1$ 为半径的圆；$x^2-y^2=0$ 表示两条相交直线 $y=x$ 与 $y=-x$
  \item[\textbf{7.}] (1) $[-5,6)$；(2) $100$ 种
  \item[\textbf{8.}] 见解析版
  \item[\textbf{9.}] 当矩形的长（落在直径上的一边）为 $\sqrt2R$、宽（垂直于直径的一边）为 $\dfrac{\sqrt2R}{2}$ 时面积最大，最大面积为 $R^2$
  \item[\textbf{10.}] $h=\dfrac{20\sqrt{52+13\sqrt3}}{13}$ 米 $=20\sqrt{\dfrac{4+\sqrt3}{13}}$ 米（约 $13.28$ 米）
  \item[\textbf{11.}] $V_{A\text{-}NBC}=108\sqrt3\ \mathrm{cm}^3$
  \item[\textbf{12.}] $2$，$6$，$18$ 或 $\dfrac29$，$-\dfrac{10}9$，$\dfrac{50}9$
  \item[\textbf{13.}] 轨迹方程为 $(x+1)^2-y^2=65$，即 $\dfrac{(x+1)^2}{65}-\dfrac{y^2}{65}=1$，它是中心在 $(-1,0)$ 的等轴双曲线
\end{itemize}
